\documentclass[11pt]{article}
\usepackage[T1]{fontenc}
\usepackage{lmodern,microtype}
\usepackage[margin=1in]{geometry}
\usepackage{amsmath,amssymb,amsthm,mathtools}
\usepackage{booktabs,array,longtable,enumitem,xcolor}
\usepackage{fancyhdr,hyperref,listings}
\hypersetup{colorlinks=true,linkcolor=blue!50!black,urlcolor=blue!50!black,citecolor=blue!50!black,pdftitle={Four rules for iPRA: substitution at witness discharge},pdfauthor={Research report prepared for Vladimir Reshetnikov}}
\pagestyle{fancy}\fancyhf{}
\fancyhead[L]{Four rules for iPRA}\fancyhead[R]{\thepage}
\setlength{\headheight}{14pt}
\setlist{itemsep=3pt,topsep=5pt}
\newtheorem{theorem}{Theorem}[section]
\newtheorem{lemma}[theorem]{Lemma}
\newtheorem{proposition}[theorem]{Proposition}
\newtheorem{corollary}[theorem]{Corollary}
\theoremstyle{definition}\newtheorem{definition}[theorem]{Definition}
\newcommand{\ipra}{\mathsf{iPRA}}
\newcommand{\K}{\mathcal K}
\newcommand{\N}{\nu}
\newcommand{\lit}{\equiv}
\newcommand{\FV}{\operatorname{FV}}
\newcommand{\subst}[2]{\N\bigl(#1[#2]\bigr)}
\newcommand{\cl}[2]{[\,#1;\,#2\,]}
\newcommand{\one}{\mathbf 1}
\newcommand{\zero}{\mathbf 0}
\newcommand{\repo}{\texttt{ConciseIPRA}}
\lstset{basicstyle=\ttfamily\scriptsize,breaklines=true,columns=fullflexible,keepspaces=true,showstringspaces=false,frame=single,rulecolor=\color{black!15},xleftmargin=3pt,xrightmargin=3pt}
\title{Four rules for intuitionistic primitive recursive arithmetic\\[5pt]
\large Substitution at witness discharge, with one computation schema}
\author{Prepared for Vladimir Reshetnikov}
\date{November 12, 2026}
\begin{document}
\maketitle
\begin{abstract}
We give a four-family presentation of full first-order intuitionistic
primitive recursive arithmetic. Terms use zero and the existing
successor-seeded recursor, with its zero-body canonicalization. Formulas
are equations and the existing paired binder $D$. The rules are one
successor calculation, clause use, paired introduction, and substitutive
witness discharge. There is no fixed reflexivity axiom, separate equality
rule, separate existential-elimination rule, or evaluation of iterator
bodies by a normalization algorithm.

The new proof is a short, noncircular recovery of equality and ordinary
witness elimination. Composed with the pinned repository's canonical
recursor and scalar bridges, it gives a whole-domain bi-interpretation
with $\ipra$, actual equality, identity numerical composites, and
equiconsistency. A matched five-family checker is reduced from 2318 to
2267 normalized bytes, from 726 to 713 significant tokens, and from 809
to 798 AST nodes. This is a concrete local simplification, not an absolute
minimum or a new repository-wide size record. The historical 516-byte
compact and 788-byte public-goal profiles are not beaten. A separate
appendix proves a constancy-only alternative to induction, while charging
its additional Boolean-equality axiom explicitly.
\end{abstract}
\setcounter{tocdepth}{1}
{\small\tableofcontents}
\newpage

\section{Result, baseline, and limits}

The question asks for a presentation simpler than the current
\repo\ project. We froze the repository at commit
\texttt{17898ea7e6a1300dfdcad71040d5daeed67916d4}. The inspected state
contains 67 filed research packages and six further ZIP archives in
\texttt{docs/incoming}. In particular, the incoming
\texttt{witness\_discharge} package already has three instruction families
and a fixed reflexivity clause, using full computation normalization.
That construction is a necessary comparison, not a new result of this
report. The existing canonical-recursors package already supplies the
successor-seeded term language used here \cite{readme,canonical,name}.

Our change is to allow a \emph{formula template} in witness discharge,
so that the discharged equality may have an arbitrary source. This
removes the dependence of the equality bootstrap on normalization of
selector packets. Explicit computation is retained, and reflexivity is
derived from it. On the canonical-recursors basis, only one explicit
computation schema is needed. Thus the complete presentation has four
families, no permanent axiom clause, and only a structural canonicalizer.

The comparison with a matched five-family implementation isolates the
fusion of equality transport and witness elimination. The stronger
claim that this presentation lies below every current complexity
frontier is \emph{not established}. It would be misleading to call four
families a global minimum when a three-family presentation is already
present, or to omit the cost of its canonicalizer. The improvement here
is the combined specification and its matched complete checker.

\paragraph{Status of the proofs.}
The new equivalence has a written proof and executable certificate checks.
It has not been newly verified in Lean or Rocq. The repository separately
formalizes the initialized-scalar/native-to-$\ipra$ bridge; its
successor-seeded canonical presentation is a written result. We use that
result with its stated hypotheses and include its complete source for
audit. We did not rebuild the upstream Lean project.

\section{The entire native syntax}

\subsection{Terms and the small canonicalization}

Start with raw numerical terms
\[
 t ::= x\mid 0\mid R(n,a;u.h).
\]
Only the body $h$ binds $u$. Both the counter $n$ and initializer $a$
are outside that binding. Define
\[
 S(a):=R(0,a;u.0),\qquad 1:=S0.
\]
The intended recursor is successor seeded:
\[
 R(0,a;u.h)=Sa,\qquad
 R(Sn,a;u.h)=h[R(n,a;u.h)/u].
\]
The first equation is made part of the syntax. Recursively normalize
all children, and replace the body of any $R$ with literal counter zero
by literal zero. Write this structural operation as $\N$. A term is
canonical exactly when every subterm $R(0,a;u.h)$ has $h\lit0$.

This operation does not unfold successor iterations, execute a body,
evaluate a numerical program, or test whether a general term denotes
zero. It preserves a nonzero outer constructor. In particular,
\[
 \N(t)\lit0\quad\Longleftrightarrow\quad t\lit0.
\]
Substitution everywhere in this calculus means capture-avoiding raw
substitution followed by $\N$. We write $F[t/v]$ below with that
canonical-substitution convention understood.

\begin{lemma}[Coherence]
For any finite term or formula $T$ and capture-free substitution $\sigma$,
\[
 \N(T[\sigma])=\N\bigl(\N(T)[\N\circ\sigma]\bigr).
\]
Canonical substitution is associative, commutes with renaming, and
introduces no new free variables beyond those in its arguments.
\end{lemma}
\begin{proof}
Use structural induction after freshening bound variables. At an
$R$-node, if the counter was zero, both sides erase the body. If a
nonzero counter becomes zero under substitution, both sides again erase
it. Otherwise both sides retain the constructor and the induction
hypotheses identify its children. Normalization cannot turn any
nonvariable constructor into literal zero. Formula binders require only
the usual lifted substitution. Associativity follows by applying the
displayed identity to successive raw substitutions.
\end{proof}

\subsection{Formulas, clauses, assumptions, and contradiction}

There is one numerical sort. Native formulas are
\[
 A ::= (s=t)\mid D_{x,y}B(x,y).
\]
The distinct variables $x,y$ are bound by $D$. Retained contexts are
finite reusable lists of clauses
\[
 \cl{A}{y.B(y)},\qquad y\notin\FV(A).
\]
Their ordinary readings are
\begin{align*}
 (s=t)^\circ&=(s=t),\\
 (D_{x,y}B)^\circ&=
 \exists x\bigl(B^\circ(x,0)\to\forall y\,B^\circ(x,y)\bigr),\\
 \cl{A}{y.B}^\circ&=A^\circ\to\forall y\,B^\circ(y).
\end{align*}
The equality on the right is actual equality. An ordinary assumption
$A$ abbreviates the constant clause $\cl{0=0}{v.A}$, with $v$ unused.
The truth of $0=0$ is \emph{derived}; it is not an initial axiom.
The designated contradiction is $0=S0$.

\section{Exactly four inference families}

Every supplied expression is canonical and well scoped. Bound names
are freshened before substitution. Clauses are reusable; no linear or
resource-sensitive reading is intended.

\paragraph{Calculation (C).}
For every counter, initializer, and body:
\[
 \frac{}{\Gamma\vdash
 R(Sn,a;u.h)=h[R(n,a;u.h)/u]}.
 \tag{C}\label{rule:c}
\]
The predecessor iterate on the right is canonicalized; at $n\lit0$
it is $Sa$. This is one nullary computation schema.

\paragraph{Use (U).}
\[
 \frac{\Gamma\vdash A}{\Gamma\vdash B(t)}
 \qquad\cl{A}{y.B(y)}\in\Gamma.
 \tag{U}\label{rule:u}
\]
The chosen clause and term $t$ are part of the certificate.

\paragraph{Introduction (I).}
\[
 \frac{\Gamma,B(t,\eta)\vdash B(t,Sz)}
      {\Gamma\vdash D_{x,y}B(x,y)},
 \qquad
 \eta=\begin{cases}
 z,&B\text{ is literally an equation and }t\lit0,\\
 0,&\text{otherwise}.
 \end{cases}
 \tag{I}\label{rule:i}
\]
Here $z$ is fresh for $\Gamma,t$ and all passive parameters of $B$.
The single local assumption is the ordinary clause just defined.
The literal test is part of the rule. In particular it is not a test
of numerical equality to zero.

\paragraph{Substitutive witness discharge (O).}
Set $\Delta_b=\Gamma,\cl{B(b,0)}{y.B(b,y)}$. Then
\[
 \frac{\Gamma\vdash D_{x,y}B\qquad
       \Delta_b\vdash s=t\qquad
       \Delta_b\vdash F(s)}
      {\Gamma\vdash F(t)}.
 \tag{O}\label{rule:o}
\]
The parameter $b$ is fresh for the outer context and the passive
parameters of $B$. The \emph{exported formula} $F(t)$ must have no free
$b$. The template $F(v)$ may be any native formula, including a constant
template; $v$ is a distinguished template variable, not the witness $b$.
The terms $s,t$ and the template may mention $b$. Requiring all of them
to be $b$-free would unnecessarily weaken the rule. All three premises
are checked in their displayed scopes.

\paragraph{No unlisted primitives.}
There is no separate hypothesis rule, reflexivity axiom, equality
transport rule, ordinary existential elimination, coverage axiom,
separation axiom, logical connective, structural inference, or general
numerical-substitution instruction. Weakening and cut are admissible
proof transformations. The sole explicit computation schema is C; the equation/progress
branch of I carries the induction strength.

\section{The equality bootstrap is noncircular}

Let
\[
 W_t:=R(1,0;u.t),\qquad J:=(W_0=0),\qquad
 T:=D_{x,y}J,
\]
where $u,x,y$ are absent from the indicated passive expressions.
All these are finite canonical expressions. Rule C gives $W_t=t$,
because the constant body ignores the predecessor iterate.

\begin{lemma}[A closed paired theorem]
The formula $T$ has a proof using C and I only.
\end{lemma}
\begin{proof}
Introduce $D_{x,y}J$ at witness $1$. The matrix is independent of both
bound variables, so the branch has goal $J$. Rule C proves $J$ while
ignoring the local assumption. In particular, this proof uses neither
ordinary-assumption use nor reflexivity.
\end{proof}

\begin{lemma}[Unrestricted formula transport]
For every native template $F(v)$,
\[
 \Gamma\vdash s=t,\quad\Gamma\vdash F(s)
 \quad\Longrightarrow\quad\Gamma\vdash F(t).
 \tag{Tr}\label{derived:tr}
\]
\end{lemma}
\begin{proof}
Use $T$ as the first premise of O. Weaken the two supplied proofs into
its opened clause context, without using the fresh witness. The exported
formula is outer scoped. This argument does not use an equality law.
\end{proof}

\begin{lemma}[Reflexivity and ordinary assumptions]
For every term $t$, $\Gamma\vdash t=t$. Consequently ordinary-assumption
use is admissible.
\end{lemma}
\begin{proof}
Calculation gives $W_t=t$. Use it as both premises of Tr with template
$F(v)=(v=t)$. Its source instance is $W_t=t$ and its target instance is
$t=t$. To use an assumption $A$, first take this construction at $t=0$;
then U on $\cl{0=0}{v.A}$ yields $A$.
\end{proof}

\begin{lemma}[Ordinary equality and ordinary witness elimination]
Symmetry, transitivity, term congruence, and the usual D-elimination
rule are admissible.
\end{lemma}
\begin{proof}
Given $s=t$, transport $s=s$ through $F(v)=(v=s)$ to obtain $t=s$.
Given $s=t$ and $t=r$, transport the second equation along $t=s$
through $F(v)=(v=r)$ to obtain $s=r$. Given $s=t$, transport the
reflexive equation $h(s)=h(s)$ through $F(v)=(h(s)=h(v))$ to get
$h(s)=h(t)$.

For ordinary D-elimination, suppose its opened branch proves an
outer-scoped formula $A$. In O choose the constant template $F(v)=A$
and use C's closed equation $J$ as the local equality premise. The
minor premise is precisely the supplied branch, and the exported
formula is $A$. This recovery does not require a naming packet or
any normalization of a selector.
\end{proof}

\paragraph{Structural proof operations.}
Weakening, exchange, contraction, and injective fresh renaming follow
by induction on finite derivations. For cut of an ordinary assumption,
replace a U node selecting its constant clause by the supplied proof;
freshen local parameters first. These arguments preserve the literal
witness test in I. We do not assume that arbitrary numerical
substitution preserves that test: a free witness can become literal
zero. The comparison below derives that admissibility indirectly.

\section{Exact equivalence with the separate-rule calculus}

Let $\K_{\rm sep}$ have the same terms, formulas, clauses, calculation,
Use and Introduction, but replace O by two rules: full Tr as displayed
above, and ordinary D-elimination. It has five families and no fixed
reflexivity axiom.

\begin{theorem}[Identity-language consequence equivalence]
For every canonical clause context $\Gamma$ and formula $A$,
\[
 \Gamma\vdash_{\K}A\quad\Longleftrightarrow\quad
 \Gamma\vdash_{\K_{\rm sep}}A.
\]
The translations are effective on finite proofs and preserve the exact
end sequent, including open numerical parameters.
\end{theorem}
\begin{proof}
Translate an O inference into Tr inside its opened witness scope,
then ordinary D-elimination. Its side condition is exactly what the
latter requires. Conversely, translate Tr using the closed theorem $T$
and O, and ordinary D-elimination using constant $F$ and local equality
$J$. All other inferences translate directly. Induct on derivations,
freshening bound eigenparameters when inserting the fixed auxiliary
proofs. No search for a numerical witness or truth value is involved.
\end{proof}

\subsection{Recovering the current canonical basis}

For an unused numerical coordinate $w$, abbreviate
\[
 q_yB(y):=D_{w,y}B(y).
\]
The following recoveries connect $\K$ to the current canonical
Eq/Q/D basis of \cite{canonical}.

First, every case introduction is admissible, including an atomic
matrix at witness zero. Let $c=W_0$. This is a nonzero \emph{syntactic}
term, and C proves $c=0$. Given a case branch at witness zero, start
instead from $B(c,0)$, transport to $B(0,0)$, insert the supplied
branch by ordinary-assumption cut, and transport its conclusion back
to $B(c,Sz)$. I at $c$ is a case introduction and has the desired
conclusion $D_{x,y}B$. No induction is used in this repair.

Next, case introduction with an unused witness gives
\[
 \Gamma,B(0)\vdash B(Sz)\quad\Longrightarrow\quad\Gamma\vdash q_yB.
\]
Ordinary D-elimination opens $q_yB$ as the clause
$\cl{B(0)}{y.B(y)}$. Together with a proof of $B(0)$, U gives $B(t)$
for any target term $t$. When $B$ is literally an equation, I at the
unused witness zero directly gives $q_yB$ from
$\Gamma,B(z)\vdash B(Sz)$. Thus contextual equation induction and
numerical cases for arbitrary native formulas are derived.

Conversely, a retained clause $\cl{B(0)}{y.B(y)}$ proves the public
formula $q_yB$: use case introduction at the unused witness $1$;
under its temporary assumption $B(0)$, use the retained clause at
$Sz$. Ordinary-assumption use is already derived. Thus a source proof
whose D-elimination branch retained a Q formula can be translated by
first deriving its dummy-D representative from the opened clause and
then cutting that proof into the translated branch. The two context
representations are related by explicit derivations, not identified
by convention.

Reflexivity, symmetry, and full transport recover the canonical
basis's crossed equation transport, including its restricted sources.
The crossed form of Q-elimination is recovered by symmetry at equation
conclusions. Conversely, in that canonical basis a clause
$\cl{A}{y.B}$ has the ordinary reading $A\to\forall y B$, represented
by its established native macros. Its derived full equality transport
and D-elimination validate O. Hence all clause sequents translate, and
ordinary formula sequents have the same consequences in both directions.

\section{The bi-interpretation with full first-order iPRA}

\subsection{What the inherited bridge actually supplies}

The predecessor construction in \cite{canonical} is made before
assuming a full numerical selector, successor cancellation, arbitrary
equality reflection, or induction on compound formulas. It supplies a
fixed canonical term $P$ with
\[
 P0=0,\qquad P(Sa)=a.
\]
Its prelogic argument supports contextual finite cases, which are
needed to avoid an illicit use of full logic during that construction.
We inherit that work rather than asserting that a counter-aware
recursor is already available.

The initialized scalar iterator is then a finite term:
\[
 \widehat I(n,a;u.h)=
 P\!\left(R\bigl(n,a;v.S(h[P(v)/u])\bigr)\right).
 \tag{Init}\label{eq:init}
\]
At zero, its computation is $P(Sa)=a$. At a successor, C exposes an
outer successor which $P$ removes. Thus
\[
 \widehat I(0,a;u.h)=a,\qquad
 \widehat I(Sn,a;u.h)=h[\widehat I(n,a;u.h)/u].
\]
The scalar bridge then gives a faithful Boolean equality test,
quantifier-free induction, pairing and projections, and a compiler for
each fixed primitive recursive description. Those implications are
explicitly proved in the included reference sources and are separately
formalized for the scalar basis in the repository \cite{native,foundations}.

\subsection{The complete logical interface is represented}

For clarity, the derived logic is not just an equational fragment.
Put $C(v,a,b)=\widehat I(v,a;u.b)$. Define the gate at equations by
\[
 G_v(s=t):=(C(v,0,s)=C(v,0,t))
\]
and extend it through $D$ by $G_v(DB)=D(G_vB)$, freshening names.
It has a provable zero instance and its successor instance is
interderivable with the gated formula. For two equation matrices set
\[
 M_v(s=t,p=q):=(C(v,s,p)=C(v,t,q)).
\]
Extend $M$ through two D roots componentwise. When only one input
has a D root, pad the other by $D_{x,y}G_y$ before descending. This
recursion decreases the maximum binder depth. It yields
$M_0(A,B)\dashv\vdash A$ and $M_{Sz}(A,B)\dashv\vdash B$.

Let $p(z)=\widehat I(z,0;u.C(u,1,0))$; its needed fixed values are
$p(1)=1$ and $p(2)=0$. The ordinary connectives and quantifiers have
the following native representatives, with fresh unused $w$:
\begin{align*}
 A\Rightarrow_*B&:=D_{w,z}M_z(A,B),\\
 A\wedge_*B&:=D_{w,z}G_zM_{p(z)}(A,B),\\
 \exists_*x\,A(x)&:=D_{x,y}G_yA(x),\\
 A\vee_*B&:=\exists_*x\,M_x(A,B),\\
 \forall_*x\,A(x)&:=A(0)\wedge_*D_{w,z}G_zA(z),\\
 \bot_*&:=(0=S0).
\end{align*}
Their usual intuitionistic introduction and elimination rules follow
from cases, the gate endpoints, clause specialization, and ordinary
D-elimination. Conjunction projections specialize at $1$ and $2$;
disjunction uses the witness as a zero/successor selector. Coverage
then follows from cases and reflexivity. From $0=S0$, selector
transport proves every equation, and external induction on formula
syntax proves every native formula. No excluded middle for quantified
formulas is introduced. The full details and both logical round trips
are retained in the included inherited logical bridge \cite{foundations}.

\begin{theorem}[Whole-domain bi-interpretation]
The ordinary realization of $\K$ and full first-order $\ipra$ are
bi-interpretable. Both domain formulas are $x=x$, interpreted equality
is actual equality, and both numerical composite isomorphisms have
graph $x=y$. Both formula composites are provably equivalent to the
original formulas. In particular the two theories are equiconsistent.
\end{theorem}
\begin{proof}
Interpret each fixed $R$ term in $\ipra$ by primitive recursion with
seed $Sa$ and its displayed update. The zero-body canonicalization
preserves its value. C is a defining equation; U is implication and
universal elimination; I uses numerical coverage in case mode and
equation induction under its assumed base in progress mode. O is
equality substitution inside an existential-elimination scope.

In the reverse direction, use the predecessor, initialized iterator
\eqref{eq:init}, scalar compiler, and the logical macros above. The
preceding section proves that their necessary rules are available in
$\K$. This supplies all defining equations and quantifier-free
induction of $\ipra$.

For numerical composites, work through each fixed finite primitive
recursive description in dependency order. After its earlier
subdescriptions have been identified, the description and its return
translation have identical seeds and updates. Equation induction on
their two trajectories proves their equality for every number in the
whole domain. For an $R$ body, pointwise body equality is lifted to
iteration by the same trajectory argument; it is not assumed as a
free intensional congruence rule. The interpretation of successor
returns $I(0,Sa;u.0)=Sa$. Hence the numerical composite maps are
the identity, not merely bijections of standard term values.

The established logical round trips identify the native $D$ reading
and its returned macro. Induction on formula construction then extends
the numerical identities to both complete formula composites. The
contradiction $0=S0$ corresponds to falsum in each direction, giving
equiconsistency by the effective proof translations.
\end{proof}

\section{The complete checker and the simplicity comparison}

\subsection{Input and trust boundary}

The selected module is \texttt{code/canonical\_checker.py}. It accepts
exactly a JSON object with \texttt{goal} and \texttt{proof}. The goal
is a closed canonical Eq/D formula. Integer zero is numerical zero,
negative integers are free eigenparameter levels, and positive integers
are de Bruijn indices. An R node has length three, an equation has
length two, and D has length one. Scope and grammar validation inspect
every child before canonicalizing; malformed content in an erased
zero-counter body is rejected.

The four instructions are \texttt{c}, \texttt{u}, \texttt{i}, and
\texttt{e}. Clause use checks its antecedent proof. Discharge checks
the D major, a local equality proof, a local template-instance proof,
and the exact exported outer-scoped goal. Ordinary assumptions are
represented by clauses with antecedent $0=0$. There is no preloaded
reflexivity clause or imported theorem database.

All replacement terms are closed with respect to loose positive
indices: their external parameters have negative levels. A replacement
term is therefore inserted unchanged beneath binders; its own internal
indices remain bound by its own R or D nodes. Template traversal still
tracks binding depth and decrements indices when a binder is opened.
This invariant is checked at the points where terms enter the kernel.

Proof and syntax recursion are structural. The only canonicalization
erases a body at a literal-zero counter. A standard coding of finite
trees therefore makes the certificate relation primitive recursive;
no uniform numerical evaluator, normalization-totality axiom, or fuel
field is required. Concrete Python recursion or memory limits can of
course reject a sufficiently large valid finite certificate.

\subsection{Measured results and fair comparisons}

All measured sources are complete standalone modules, including JSON
input, validation, binding operations, inference checks, and final-goal
comparison. Actual module docstrings are removed; the AST is serialized
by \texttt{ast.unparse} with a final LF. Token counts omit layout and
comments. Gzip-9 with timestamp zero is informational.

\begin{center}\small
\begin{tabular}{lrrrr}
\toprule
Complete module & Bytes & Tokens & AST & Gzip\\
\midrule
Matched separate five-family control & 2318 & 726 & 809 & 889\\
Selected four-family calculus & 2267 & 713 & 798 & 882\\
\bottomrule
\end{tabular}
\end{center}
The reduction is 51 bytes (2.20\%), 13 significant tokens (1.79\%),
and 11 AST nodes (1.36\%). These are modest savings. The control and
proposal share the same grammar, canonicalizer, binding representation,
assumption representation, names, and I rule; only equality transport
and witness elimination are separated or fused.

The incoming full-normalization Name checker is 2273 normalized bytes,
732 tokens, and 813 AST nodes as delivered. The numerical comparison
2267 versus 2273 changes several design decisions and is \emph{not}
an isolated measure of our mathematical change. The current compact
516-byte and validated-public-goal 788-byte profiles are much smaller
and use different certificate contracts. Our result is not a new
corpus-wide minimum.

\subsection{Executed checks}

The canonical suite passed 14 positive and 18 adversarial cases, 14
exact-goal translations through the matched control and back, and 64
standalone checks across ordinary and optimized Python. The raw companion
passed 14 positives, 16 rejections, two open sequents, 60 command-line
checks, 14 control round trips, and five translations from the frozen
raw CE5 checker. There were no unexpected outcomes.

The archive includes executable positive and adversarial certificates,
an independently generated raw-scalar companion, and the matched
canonical control. The final evidence file records the exact runs.
The intended coverage is computation, reflexivity derived without an
axiom, use of a retained induction hypothesis, compound-template
transport, nested binding, ordinary witness elimination, and rejection
of witness escape and ill-scoped zero-body contents. Finite tests
validate these implementations; the general soundness, equivalence,
and bi-interpretation claims rest on the proofs above and the stated
inherited bridges.

\section{A boundary result: unary congruence is insufficient}

One tempting further simplification is to replace binary equality
transport by the unary rule $s=t\Rightarrow h(s)=h(t)$. The following
small counterinterpretation shows why that does not work for the
initialized-scalar raw or zero-normal language, even when reflexivity,
symmetry, computation, and equation induction are retained.

\begin{proposition}
In that reduced Eq/Q/D calculus, transitivity of equality is not
derivable. This remains true when $I(0,a;u.h)$ is definitionally $a$.
\end{proposition}
\begin{proof}
Use $V=\{0,m,2\}$, identified for distance purposes with
$\{0,1,2\}$. Interpret zero as 0, successor as the constant 2, and
\[
 I(n,a;u.h)=\begin{cases}a,&n=0,\\m,&n\ne0.\end{cases}
\]
The body is ignored. Interpret the syntactic equation atom by the
reflexive symmetric relation $R(s,t)$ meaning that their numerical
distance is at most one. This is a truth assignment to judgments,
not a structure using actual first-order equality. The contradiction
$0=S0$ is false.

Zero computation holds as identity. A successor iterator denotes
$m$, which is related to every possible value of its right-hand side,
so successor computation holds as a judgment. The numerical operation
$F(n,a)$ above preserves $R$ simultaneously in both arguments: when
exactly one counter is zero the outputs are $a$ and $m$; when both
are zero use $R(a,a')$; otherwise both outputs are $m$. Thus every
term template preserves $R$, validating unary congruence.

Each term function is either constantly $m$, or equals a payload
$p\in\{0,2,\text{one free variable}\}$ when all variables in a finite
gate set are zero, and equals $m$ otherwise. Structural induction proves
this classification. Fixing all but one variable, the possible value
triples are
\[
 (0,0,0),(m,m,m),(2,2,2),(0,m,2),(0,m,m),(2,m,m).
\]
Therefore an atom's one-variable truth set is empty, $\{0,m\}$,
$\{m,2\}$, or all of $V$.

Give Q and D their displayed Boolean readings. Q applied to an atom
is true, false, or a finite disjunction of conditions that a free
variable is nonzero. Indeed, if its newly bound variable is in a
term's gate, the atom is true at both nonzero values; if it is only
in a payload, the only obstruction to closure occurs when a finite
set of other variables are all zero. Existential binding for D leaves
such a disjunction alone or makes it true. A further Q or D is true.
It follows that every native formula satisfies
$B(0)\land B(2)\Rightarrow B(m)$ under fixed remaining parameters.

This property validates case introduction because every successor is
2. It also validates equation progress: from its base at 0 and its
step at 0 obtain the value at 2, then the value at $m$. QE and DE are
sound under their ordinary Boolean readings. Hypotheses and the other
retained rules preserve truth as well.

Assign $x=0,y=m,z=2$. Both $x=y$ and $y=z$ are true, while $x=z$
is false. Thus the sequent expressing transitivity cannot be derived.
\end{proof}

Full successor-computation normalization is outside this obstruction:
it would identify values which this interpretation only relates.
Nor does the result exclude every different arithmetic presentation.
It excludes this specific deletion in a precisely stated neighborhood
of the current scalar core.

\section{Further research questions}

\begin{enumerate}
\item Can substitutive witness discharge be formalized as a small
comparison module on top of the existing Lean clause development?
The closed dummy-D proof makes its two translations especially direct.
\item Can its template be restricted to a term template, with compound
transport recovered, while preserving a net saving in the complete
checker? The permitted witness occurrences must be tracked carefully.
\item Can a guard independent of the supplied witness replace I's
literal-zero test without undoing the source-size reduction?
\item Can the constancy-only presentation in the next appendix recover
its Boolean equality split from a simpler arithmetic basis? The split
must not be silently assumed during that bootstrap.
\item Can a common public grammar and malformed-input contract be fixed
for the compact and strict interfaces? Only then would a lower bound
or a global checker-size claim have an unambiguous target.
\end{enumerate}

\appendix
% Standalone appendix fragment.  Requires amsmath, amssymb, amsthm.
% The only ambient notation is ordinary intuitionistic first-order logic
% with actual equality.  This is a proved alternative induction interface,
% not a claim of a globally smaller complete checker.
\section{A constant-recurrence rule in place of induction}
\label{sec:constancy-with-split}

This section gives a separate, complete arithmetic presentation.  Its
primitive inference rule has neither an induction hypothesis nor an
arbitrary update function: it says that a term which is unchanged by
successor has its initial value everywhere.  The price is an explicit
Boolean equality operation.  That price is part of the presentation and
must be counted in any simplicity comparison.

\subsection{The complete arithmetic base}

Use ordinary intuitionistic first-order logic with actual equality, on
one numerical sort.  Terms have the grammar
\[
 t ::= x\mid 0\mid S(t)\mid \delta(t,t)
          \mid R(n,a;i,u.h).
\]
The last constructor binds $i,u$ in $h$ only.  Its body is finite earlier
term syntax; it is not a variable ranging over numerical functions.
All substitutions are capture avoiding.  Put $1=S0$ and
$s\ne t=(s=t\to\bot)$.  The nonlogical axioms are the universal closures
of the following four displayed schemata:
\begin{align}
 R(0,a;i,u.h)&=a, \tag{R0}\label{eq:const-r0}\\
 R(Sn,a;i,u.h)&=h[n/i,R(n,a;i,u.h)/u],
       \tag{RS}\label{eq:const-rs}\\
 0=1&\longrightarrow\bot, \tag{Sep}\label{eq:const-sep}\\
 \bigl(s=t\land\delta(s,t)=0\bigr)&\ \lor\
 \bigl(s\ne t\land\delta(s,t)=1\bigr).
       \tag{Split}\label{eq:const-split}
\end{align}
The substitution in \textup{(RS)} is simultaneous.  No coverage axiom,
successor-injectivity axiom, induction schema, or equality-reflection
theorem is silently included.

The one additional inference schema is
\[
 \frac{\Gamma\vdash f(Sz)=f(z)}
      {\Gamma\vdash f(t)=f(0)}\quad(\mathrm{Const}),
 \tag{Const}\label{eq:const-rule}
\]
where $f(x)$ is a fixed numerical term template, $z$ is fresh for
$\Gamma,t$ and the passive parameters of $f$, and the displayed premise
contains no new numerical or equational assumption.  The context
$\Gamma$ may contain arbitrary first-order formulas.

The underlying logic can be taken to be standard intuitionistic natural
deduction: assumption; introduction and elimination for
$\land,\lor,\to,\forall,\exists$; falsum elimination; equality reflexivity
and substitution; and the usual eigenvariable conditions for
$\forall$ introduction and $\exists$ elimination.  Thus this presentation
does not use a classical inference rule or a choice principle.

\subsection{Elementary computations require no induction}

Define
\[
 C(v,a,b)=R(v,a;i,u.b),\qquad
 a+b=R(b,a;i,u.Su),\qquad
 P(v)=R(v,0;i,u.i).
\]
Their displayed equations are direct instances of the computation axioms:
\begin{align*}
 C(0,a,b)&=a,& C(Sv,a,b)&=b,\\
 a+0&=a,& a+Sb&=S(a+b),\\
 P0&=0,& P(Sv)&=v.
\end{align*}
In particular $a+1=Sa$ and successor is injective.

Every successor is different from zero.  Indeed, if $Sv=0$, applying
$w\mapsto C(w,0,1)$ gives $1=0$, contradicting \textup{(Sep)}.
This argument uses no case distinction on an arbitrary number.

The split axiom, together with $0\ne1$, gives
\[
 \delta(s,t)=0\lor\delta(s,t)=1,
 \qquad \delta(s,t)=0\ \longleftrightarrow\ s=t.
 \tag{E}\label{eq:const-eqreflection}
\]
For the reverse direction of the equivalence, eliminate the two
disjuncts of \textup{(Split)} under $s=t$; its second disjunct contradicts
that assumption.  For the forward direction, its second disjunct would
give $0=1$.  These are intuitionistic derivations.

\subsection{Boolean failure terms for quantifier-free formulas}

For every quantifier-free formula $A$, recursively define a term $q_A$,
with value zero meaning that $A$ holds:
\begin{align*}
 q_{s=t}&=\delta(s,t), &q_\bot&=1,\\
 q_{A\land B}&=C(q_A,q_B,1),
 &q_{A\lor B}&=C(q_A,0,q_B),\\
 q_{A\to B}&=C(q_A,q_B,0).
\end{align*}
Negation is the abbreviation $A\to\bot$.  Structural induction on the
formula, using the two computations for $C$, proves
\[
 q_A=0\lor q_A=1,
 \qquad q_A=0\ \longleftrightarrow\ A.
 \tag{B}\label{eq:const-boolean}
\]
The induction here is a metatheoretic induction on a finite formula, not
an arithmetic induction inside the proposed theory.  In the implication
case, split on $q_A$: when it is zero, $A$ is available and $A\to B$ is
equivalent to $B$; when it is one, $\neg A$ follows and the implication
is true.  Only the decidability already proved for the quantifier-free
antecedent is used.  No claim is made for formulas with quantifiers.

\subsection{An accumulator turns progress into constancy}

Suppose $A(n)$ is quantifier free and that the proposed theory has
derivations
\[
 \Gamma\vdash A(0),\qquad
 \Gamma,A(z)\vdash A(Sz),
\]
with the usual fresh-variable condition.  Write $q(n)=q_{A(n)}$.
Define the following finite terms:
\begin{align*}
 e(n)&=C\bigl(q(n),0,C(q(Sn),1,0)\bigr),\\
 a(n)&=R(n,0;i,u.u+e(i)),\\
 F(n)&=a(n)+q(n).
\end{align*}
Here $e$ and the addition operation are expanded into their finite term
abbreviations before the outer $R$ is formed.  In particular, the
definition of $a$ does not presuppose a new function symbol or a new
recursion principle.

The computations give
\[
 a(0)=0,\qquad a(Sn)=a(n)+e(n).
\]
The three possible pairs of Boolean values of $q(n),q(Sn)$ give the
following equalities without an induction hypothesis:
\[
\begin{array}{c|c|c|c}
 q(n)&q(Sn)&e(n)&F(Sn)=F(n)\\ \hline
 0&0&0&a(n)=a(n)\\
 1&1&0&Sa(n)=Sa(n)\\
 1&0&1&Sa(n)=Sa(n)
\end{array}
\]
The remaining pair $(0,1)$ is impossible: $q(n)=0$ gives $A(n)$;
the retained step derivation gives $A(Sn)$; and
\eqref{eq:const-boolean} then gives $q(Sn)=0$, contradicting $0\ne1$.
Thus, by finite disjunction elimination under $\Gamma$,
\[
                       \Gamma\vdash F(Sz)=F(z).
\]
Apply \textup{(Const)} to obtain $F(t)=F(0)$.  The proved base $A(0)$
gives $q(0)=0$, so $F(0)=0$.  If $q(t)=1$, then
\[
                 0=F(t)=a(t)+1=Sa(t),
\]
contrary to successor separation.  The Boolean disjunction for $q(t)$
therefore gives $q(t)=0$, and hence $A(t)$.

\begin{theorem}[Recovery of contextual quantifier-free induction]
The four axiom schemata \textup{(R0)}, \textup{(RS)}, \textup{(Sep)},
\textup{(Split)} and the rule \textup{(Const)} derive
\[
 \frac{\Gamma\vdash A(0)\qquad\Gamma,A(z)\vdash A(Sz)}
      {\Gamma\vdash A(t)}
\]
for every quantifier-free $A$, under arbitrary retained first-order
contexts.  In particular they derive contextual induction for every
literal equation.
\end{theorem}

No commutativity or associativity of addition was used in this proof.
The only addition identities used are $a+0=a$ and $a+1=Sa$.
This is why the arithmetic bootstrap does not conceal an induction
argument in an algebraic rearrangement.

\subsection{Coverage, primitive recursion, and both interpretations}

Apply the recovered rule to the quantifier-free formula
\[
                  A(n)\equiv n=0\lor n=S(Pn).
\]
Its zero case is immediate.  At a successor, $P(Sn)=n$ gives the second
disjunct without an induction hypothesis.  Thus numerical coverage is
derived.  Successor injectivity was already obtained from $P(Sn)=n$.

Every usual primitive recursive definition compiles directly into this
counter-aware language: zero, successor, projections, and composition use
the corresponding term operations, and
\[
 f(0,\vec x)=g(\vec x),\qquad
 f(Sn,\vec x)=h(n,f(n,\vec x),\vec x)
\]
is represented by
\[
 R\bigl(n,g(\vec x);i,u.h(i,u,\vec x)\bigr),
\]
after recursively compiling $g,h$.  The computation equations are
literal instances of \textup{(R0)} and \textup{(RS)}.  Together with
successor separation, injectivity, and recovered quantifier-free
induction, this gives an interpretation of full first-order iPRA.

Conversely, each fixed finite $R$ term denotes a primitive recursive
function in iPRA by ordinary primitive recursion on its counter.  Its
two defining equations hold by construction.  Interpret $\delta$ by a
Boolean equality characteristic.  For an explicit construction, define
\[
                         d(a,b)=R(b,a;i,u.Pu).
\]
The subtraction equations and equation induction give
\[
 d(a,0)=a,\quad d(a,Sb)=P(d(a,b)),\quad
 d(Sa,Sb)=d(a,b),\quad d(Sn,n)=1.
\]
For the third identity, induct on $b$ at fixed $a$; the base is
$P(Sa)=a$, and the step applies $P$ to the preceding equality.
For the fourth, induct on $n$ using the third identity.  Elementary
equation inductions on the second argument of addition also give
$0+b=b$ and $Sa+b=S(a+b)$.

Quantifier-free induction on $b$, with $a$ passive, proves
\[
           d(a,b)\ne0\ \longrightarrow\ a=b+d(a,b).
 \tag{Rec}\label{eq:const-reconstruct}
\]
The base uses $0+a=a$.  For the step, put $r=d(a,b)$ and suppose
$Pr=d(a,Sb)\ne0$.  Coverage on $r$ gives $r=S(Pr)$: the zero case
contradicts $Pr\ne0$, and the successor case uses the equation for $P$.
In particular $r\ne0$, so the preceding induction hypothesis gives
\[
              a=b+r=b+S(Pr)=Sb+Pr.
\]
This establishes the new instance of \eqref{eq:const-reconstruct}.

Now define
\[
 Z(a)=C(a,1,0),\qquad U(a)=C(a,0,Z(Pa)),\qquad
 E(a,b)=C\bigl(b,Z(a),U(d(a,Pb))\bigr).
\]
Coverage and successor separation give $Z(a)\in\{0,1\}$ and
$Z(a)=1\leftrightarrow a=0$.  They similarly give
$U(a)\in\{0,1\}$ and $U(a)=1\leftrightarrow a=1$.
If $b=0$, $E(a,b)=1$ is equivalent to $a=0$.
If $b=Sn$, $E(a,b)=1$ is equivalent to $d(a,n)=1$.
The latter condition implies $a=n+1=Sn=b$ by
\eqref{eq:const-reconstruct}; conversely $a=Sn$ gives $d(Sn,n)=1$.
Thus $E(a,b)$ is Boolean and equals one exactly when $a=b$.
The fixed primitive recursive term
\[
                         \delta_*(a,b)=Z(E(a,b))
\]
is therefore Boolean and equals zero exactly when $a=b$.  This
derivation uses only quantifier-free inductions after they have already
been recovered; it is not part of the induction-free bootstrap.

Consequently the chosen equality characteristic satisfies
\textup{(Split)}.  The rule \textup{(Const)} is iPRA-admissible:
under its recurrence premise, equation induction on
$f(n)=f(0)$ gives the conclusion.  All domains remain the whole
numerical domain and equality is interpreted as actual equality.

For completeness, the interpretation composites require more than
agreement in the standard natural numbers.  In the proposed theory,
recovered equation induction proves uniqueness of every fixed primitive
recursive defining recursion, under arbitrary parameters.  By structural
induction on term syntax, it therefore proves that the round trip of
every $R$ term has the same value as the original.  The same argument in
iPRA treats the round trip of each standard primitive recursive function
symbol.  For the additional primitive $\delta$, both $\delta$ and the
translated $\delta_*$ satisfy \textup{(Split)} in the proposed theory.
Case analysis on the first split shows their equality: if $a=b$, both
are zero; if $a\ne b$, both are one.  Hence
\[
                         \delta(a,b)=\delta_*(a,b)
\]
is provable there.  Equality congruence and structural induction on
formulas now give both formula round trips.  The numerical isomorphisms
have graph $x=y$.

\begin{theorem}[Exact scope of the result]
The presentation in this section is whole-domain bi-interpretable with
full first-order iPRA, with both composite numerical interpretations
provably the identity.  It is therefore equiconsistent with iPRA.
The universal statements are written mathematical proofs, not a new
proof-assistant formalization.
\end{theorem}

\subsection{What this does and does not simplify}

The primitive arithmetic inference now concerns a single numerical term
and an unconditional identity-update recurrence.  It has no discharged
induction hypothesis, no Boolean accumulator in its rule text, and no
arbitrary update template.  The accumulator is an untrusted derived
proof construction.

The separate equality split is essential to the proof given here.
Replacing it by an assertion that an internally defined equality test
is faithful, without proving that assertion in the induction-free base,
would be circular.  Likewise, for a tail-recursive iterator,
$d(a,Sb)=d(Pa,b)$ does not by itself imply
$d(a,Sb)=P(d(a,b))$: the missing commutation is a fusion theorem.

Compared with the repository's scalar basis, this presentation adds a
binary term constructor, a stronger primitive equality split in place
of primitive coverage, and a second local binder for the recursor.
It removes arbitrary induction from the primitive inference interface.
No claim is made that those changes reduce the complete checker-size
metric, surpass the existing compact profiles, or establish an absolute
minimum.  The unanswered next question is whether the equality split
can be derived from a sufficiently small recursion base and
\textup{(Const)} alone.


\section{Reproduction and artifact map}

Run \texttt{python3 reproduce.py} in the extracted directory. It
replays the included validation scripts and regenerates complete-module
metrics. Run \texttt{python3 build\_pdf.py} to rebuild this PDF with
pdfLaTeX; no shell escape is needed. The exact selected checker is
included below and as an executable file. Supporting source snapshots
retain the repository's MIT No Attribution license.

\begin{description}[style=nextline]
\item[\texttt{report.tex}, \texttt{report.pdf}]
The new construction, proofs, explicit limits, and the constancy appendix.
\item[\texttt{code/}]
Complete selected and matched-control checkers.
\item[\texttt{validation/}, \texttt{examples/}, \texttt{evidence/}]
Untrusted certificate producers, executable checks, measured sizes,
and results from the delivered run.
\item[\texttt{reference/}]
Pinned current README, canonical-recursors proof source, logical and
ordinary bridges, source metadata, and license. These are attributed
prerequisites; their historical novelty claims are not claims of this report.
\item[\texttt{CLAIMS.json}, \texttt{FORMAT.txt}, \texttt{README.txt}]
Machine-readable scope, the exact certificate contract, and entry points.
\end{description}

\clearpage
\section{Complete selected acceptance module}
\begin{lstlisting}[language=Python]
"""Canonical successor-seeded iPRA with substitutive witness discharge.

Terms: integer 0; negative free levels; positive bound indices; [n,a,h].
R(0,a;u.h) is canonically [0,a,0], and S(a) abbreviates [0,a,0].
Formulas: [s,t] or [B], with D binding x at index 2 and y at index 1.
Four instructions: c (successor calculation), u (clause use),
i (paired introduction), e (substitutive witness discharge).
No reflexivity axiom and no execution of iterator bodies by normalization.
Input: exactly {'goal': closed_formula, 'proof': finite_proof_tree}.
"""
import json
import sys


def need(condition):
    if not condition:
        raise ValueError('invalid certificate')


def walk(a, f, d=0, form=None):
    if type(a) is int:
        need(not form)
        return f(a, d)
    need(type(a) is list)
    size = len(a)
    need(form is None or size in ((1, 2) if form else (3,)))
    a = [walk(b, f, d + (2 if size == 1 else int(size == 3 and i == 2)), size == 1)
         for i, b in enumerate(a)]
    if size == 3 and a[0] == 0:
        a[2] = 0
    return a


def syntax(a, n, form=False, depth=0):
    def variable(i, d):
        need(-n <= i <= d)
        return i
    need(walk(a, variable, depth, form) == a)


def sub(a, t, depth=0):
    return walk(a, lambda i, d: t if i == d + 1 else i - (i > d + 1), depth)


H = [0, 0]


def prove(a, p, g, n):
    syntax(a, n, True)
    match p:
        case ['c']:
            need(len(a) == 2)
            counter, seed, body = a[0]
            need(type(counter) is list and counter[0] == 0)
            prev = [counter[1], seed, body if counter[1] else 0]
            need(a[1] == sub(body, prev))
        case ['u', j, t, premise]:
            need(type(j) is int and 0 <= j < len(g))
            syntax(t, n)
            base, body = g[j]
            need(a == sub(body, t))
            prove(base, premise, g, n)
        case ['i', t, branch]:
            need(len(a) == 1)
            syntax(t, n)
            body = sub(a[0], t, 1)
            hyp = sub(body, -n - 1 if len(body) == 2 and t == 0 else 0)
            prove(sub(body, [0, -n - 1, 0]), branch, g + [(H, hyp)], n + 1)
        case ['e', body, fml, s, t, major, eq, branch]:
            prove([body], major, g, n)
            body = sub(body, -n - 1, 1)
            local = g + [(sub(body, 0), body)]
            syntax(fml, n + 1, True, 1)
            prove([s, t], eq, local, n + 1)
            prove(sub(fml, s), branch, local, n + 1)
            need(a == sub(fml, t))
        case _:
            raise ValueError('invalid proof')


def check(data):
    need(type(data) is dict and set(data) == {'goal', 'proof'})
    prove(data['goal'], data['proof'], [], 0)


if __name__ == '__main__':
    try:
        check(json.load(sys.stdin))
    except Exception:
        print('REJECT')
        sys.exit(1)
    print('ACCEPT')
\end{lstlisting}

\begin{thebibliography}{9}
\raggedright
\bibitem{readme} Vladimir Reshetnikov, \emph{ProveIt: Concise presentations
of intuitionistic primitive recursive arithmetic}. Pinned commit
\texttt{17898ea7e6a}, current README and incoming archive inventory.
\url{https://github.com/VladimirReshetnikov/ProveIt/tree/17898ea7e6a1300dfdcad71040d5daeed67916d4/Logic/PrimitiveRecursiveArithmetic/ConciseIPRA}.
\bibitem{canonical} ProveIt research package, \emph{Canonical recursors:
iPRA with one explicit computation schema}, \texttt{Research/canonical\_recursors/}.
The complete pinned report source is included under \texttt{reference/}.
\bibitem{name} ProveIt incoming package, \emph{iPRA with equality at witness
discharge}, \texttt{docs/incoming/ipra\_witness\_discharge\_2026-11-11.zip}.
Three families, fixed reflexivity clause, and full normalization.
\bibitem{native} ProveIt, \texttt{Lean/ConciseIPRA/NativeBiInterpretation.lean},
\texttt{NativeClausal.lean}, and \texttt{NativeUniformClauses.lean}, at the
pinned commit. Deep embeddings and the existing scalar/native comparison.
\bibitem{foundations} ProveIt, \texttt{Research/certified\_zero/reference/}:
\texttt{inherited\_choice\_basis.tex}, \texttt{inherited\_ordinary\_bridge.tex},
\texttt{scalar\_predecessor.tex}, \texttt{equality\_reflection.tex}, and
\texttt{primitive\_recursion.tex}. Attributed source files are included.
\end{thebibliography}
\end{document}
